\documentclass[11pt]{article}
\usepackage{fontspec}
\setmainfont{Times New Roman}
\setsansfont{Arial}
\setmonofont{Consolas}
\usepackage{amsmath,amssymb,mathtools}
\usepackage{geometry}
\usepackage{microtype}
\geometry{a4paper,margin=2.35cm}
\setlength{\parindent}{1.25em}
\setlength{\parskip}{0pt}
\makeatletter
\renewcommand\footnotesize{\@setfontsize\footnotesize{7.7}{8.7}\selectfont}
\makeatother
\emergencystretch=100em
\allowdisplaybreaks
\sloppy
\hbadness=10000
\vbadness=10000
\hfuzz=2pt
\vfuzz=10pt
\raggedbottom

\begin{document}

% Unit: Paper31 Section05 entry 1 v001. Direct Interslavic Latin translation from the German final-audited source slice, source lines 307--317 / segments P31-S0065--P31-S0068, expanded against printed scan/OCR witness page 95. Footnote counter continues after Paper31 Section04 entry 8.
\setcounter{footnote}{26}

\subsection*{§5. Direktne sumy klasov}

\textbf{1. \emph{Direktna suma klasa idealov ili klasa predstavjenj.}} Ako ideal $\mathfrak a$ jest direktna suma, $\mathfrak a=\mathfrak b+\mathfrak c$, togda iz izomorfizma slěduje, že vsaky ideal togo klasa takože stava direktna suma,
$\bar{\mathfrak a}=\bar{\mathfrak b}+\bar{\mathfrak c}$. Pri tom $\bar{\mathfrak b}$ jest odpovědajuče izomorfno k $\mathfrak b$, a $\bar{\mathfrak c}$ k $\mathfrak c$, ako idealy razmatrivajut se kako $\mathfrak R$-moduly; zato možno govoriti o direktnoj sumě klasa idealov i o komponentnyh klasah.

Ako kolco $R_{\mathfrak a}$ někojego klasa predstavjenj jest direktna suma podkolc, ktore Jednočasno sut idealy v $R_{\mathfrak a}$, togda zaradi izomorfizma to samo važi za vsako kolco $R_{\mathfrak b}$ togo klasa predstavjenj, pri čem komponenty stavajut ekvivalentne kolca. Zato možno govoriti o direktnoj sumě klasa predstavjenj i o komponentnyh klasah.

Direktnoj sumě klasa idealov odgovarja direktna suma klasa predstavjenj; v tom smyslu budemo govoriti o direktnoj sumě klasa.\footnote{Pitanje, v kakvom obsegu vse direktne sumy klasov predstavjenj mogut byti stvorjene direktnymi sumami klasov idealov, tuto znovu ne razsmatrja se.} Naj $\beta_1,\ldots,\beta_m$ bude baza $\mathfrak b$, a $\gamma_1,\ldots,\gamma_n$ baza $\mathfrak c$; poněže suma jest direktna, $\beta$ i $\gamma$ zajedno tvorjat linearno nezavisnu bazu $\mathfrak a$. Ako $R_{\mathfrak a}$ jest matrično kolco stvorjeno posrědstvom toj bazy, togda matrica iz $R_{\mathfrak a}$, prirěđena libovoljnomu elementu $\delta$ iz $\mathfrak R$, imaje formu
\[
        \begin{pmatrix}D^{\mathfrak b}&0\\0&D^{\mathfrak c}\end{pmatrix}.
\]
Zato opredělimo $R^{\mathfrak b}$ kako sistem vsih matric
\[
        \begin{pmatrix}D^{\mathfrak b}&0\\0&0\end{pmatrix}
\]
i odpovědajuče opredělimo $R^{\mathfrak c}$. Trěba pokazati, že $R^{\mathfrak b}$ i $R^{\mathfrak c}$ tvorjat podkolca v $R_{\mathfrak a}$, i že $R_{\mathfrak a}$ jest jih direktna suma; Jednočasno $R^{\mathfrak b}$ i $R^{\mathfrak c}$ stavajut izomorfne k kolcam $R_{\mathfrak b}$ i $R_{\mathfrak c}$, stvorjenym odpovědajuče iz $\mathfrak b$ i $\mathfrak c$.

Bo črez homomorfizm $R^{\mathfrak b}$ odgovarja vsim takim elementam $\delta$ iz $\mathfrak R$, za ktore $\delta\mathfrak c$ izčezaje, to jest ktore prinadležet k idealnomu kvocientu nulovogo ideala črez $\mathfrak c$, $(0):\mathfrak c$; zaradi homomorfizma $R^{\mathfrak b}$ tym samym jest idealom v $R_{\mathfrak a}$, i to samo važi za $R^{\mathfrak c}$. Kolco $R_{\mathfrak a}$ jest suma tych idealov, i imeno direktna suma, bo produkt $R^{\mathfrak b}R^{\mathfrak c}$ dvuh kolc izčezaje. Nakonec, prirěđenje, ktore matrici $D^{\mathfrak b}$ prirěđuje matricu iz $R^{\mathfrak b}$, sodržajuću $D^{\mathfrak b}$ i nuly, daje izomorfizm medžu $R^{\mathfrak b}$ i $R_{\mathfrak b}$; v tom smyslu komponentnym klasam $(\mathfrak b)$ i $(\mathfrak c)$ klasa idealov $\mathfrak a$ odgovarjajut komponentne klasy kolca $R_{\mathfrak a}$.

\end{document}
